\documentclass[12pt]{article}
\usepackage{amsmath, amssymb}
\usepackage[margin=1in]{geometry}
\setlength{\parskip}{6pt}

\title{QUBO Formulation: Traveling Salesperson Problem}
\date{}

\begin{document}
\maketitle

\subsection*{Traveling Salesperson Problem}

\paragraph{Problem Statement.}
Given a complete directed graph $G = (V, E)$ with $|V| = N$ cities and a non-negative distance matrix $D = (d_{ij})_{i,j \in V}$, find a Hamiltonian cycle that visits every city exactly once and returns to the starting city, such that the total travel distance is minimized.

\paragraph{Decision Variables.}
For every ordered pair of distinct cities $(i, j)$ with $i, j \in V$ and $i \neq j$, we introduce a binary decision variable $x_{ij} \in \{0, 1\}$. The variable $x_{ij}$ equals $1$ if the directed edge $(i, j)$ is included in the tour, and $0$ otherwise.

\paragraph{QUBO Derivation.}
We construct a Quadratic Unconstrained Binary Optimization (QUBO) Hamiltonian $H(\mathbf{x})$ by combining the objective function with quadratic penalty terms that enforce the problem constraints. Let $\lambda > 0$ be a penalty weight parameter.

\textbf{Step 1 -- Objective Function.}
The original minimization objective is the total travel distance:
\begin{align}
    O(\mathbf{x}) = \sum_{i \in V} \sum_{j \in V, j \neq i} d_{ij} x_{ij}.
\end{align}

\textbf{Step 2 -- Constraint Formulation.}
The problem requires three sets of constraints:
\begin{align}
    \text{(Out-degree)} \quad & \sum_{j \in V, j \neq i} x_{ij} = 1, \quad \forall i \in V, \\
    \text{(In-degree)} \quad & \sum_{i \in V, i \neq j} x_{ij} = 1, \quad \forall j \in V, \\
    \text{(Subtour elimination)} \quad & \sum_{i \in S} \sum_{j \in S, j \neq i} x_{ij} \le |S| - 1, \quad \forall S \subset V, \ 2 \le |S| \le N-1.
\end{align}

\textbf{Step 3 -- Penalty Term Expansion.}
Each equality constraint is converted to a quadratic penalty. For the out-degree constraint associated with city $i$, we write:
\begin{align}
    P_i^{\text{out}}(\mathbf{x}) = \lambda \left( \sum_{j \neq i} x_{ij} - 1 \right)^2.
\end{align}
Expanding the square and applying the idempotent property $x_{ij}^2 = x_{ij}$ for binary variables yields:
\begin{align}
    \left( \sum_{j \neq i} x_{ij} - 1 \right)^2 
    &= \sum_{j \neq i} x_{ij}^2 + \sum_{j \neq k, j,k \neq i} x_{ij}x_{ik} + 1 - 2\sum_{j \neq i} x_{ij} \\
    &= 1 - \sum_{j \neq i} x_{ij} + \sum_{j \neq k, j,k \neq i} x_{ij}x_{ik}.
\end{align}
Summing over all $i \in V$, the total out-degree penalty contribution is:
\begin{align}
    P^{\text{out}}(\mathbf{x}) = \lambda \sum_{i \in V} \left( 1 - \sum_{j \neq i} x_{ij} + \sum_{j \neq k, j,k \neq i} x_{ij}x_{ik} \right).
\end{align}
Analogously, the in-degree penalty for city $j$ expands to:
\begin{align}
    P^{\text{in}}(\mathbf{x}) = \lambda \sum_{j \in V} \left( 1 - \sum_{i \neq j} x_{ij} + \sum_{i \neq k, i,k \neq j} x_{ij}x_{kj} \right).
\end{align}
The subtour elimination constraints are enforced via a similar quadratic penalty over all valid subsets $S$:
\begin{align}
    P^{\text{sub}}(\mathbf{x}) = \lambda \sum_{\substack{S \subset V \\ 2 \le |S| \le N-1}} \left( \sum_{i \in S} \sum_{j \in S, j \neq i} x_{ij} - (|S|-1) \right)^2.
\end{align}

\textbf{Step 4 -- Combined Hamiltonian.}
The full QUBO Hamiltonian is the sum of the objective and all penalty terms:
\begin{align}
    H(\mathbf{x}) = O(\mathbf{x}) + P^{\text{out}}(\mathbf{x}) + P^{\text{in}}(\mathbf{x}) + P^{\text{sub}}(\mathbf{x}).
\end{align}

\textbf{Step 5 -- Q Matrix Entries and Offset.}
Collecting linear, quadratic, and constant terms from the expanded form above, the QUBO parameters are given in closed form as:
\begin{align}
    Q_{ij,ij} &= d_{ij} - \lambda, \quad \forall i \neq j, \\
    Q_{ij,ik} &= 2\lambda, \quad \forall i \in V, \ j \neq k, \\
    Q_{ji,jk} &= 2\lambda, \quad \forall j \in V, \ i \neq k, \\
    Q_{ij,kl} &= 2\lambda, \quad \text{for pairs appearing in subtour subsets } S, \\
    c &= 2N\lambda + \lambda \sum_{\substack{S \subset V \\ 2 \le |S| \le N-1}} (|S|-1)^2.
\end{align}
The Hamiltonian takes the standard QUBO form $H(\mathbf{x}) = \sum_{i} Q_{ii} x_i + \sum_{i < j} Q_{ij} x_i x_j + c$.

\paragraph{Penalty Weight Justification.}
The penalty weight $\lambda$ must be chosen sufficiently large to ensure that any constraint violation is strictly more costly than the maximum possible reduction in the objective function. Let $D_{\max} = \max_{i,j \in V} d_{ij}$. If a single degree constraint is violated, the penalty increases by at least $\lambda$. The maximum possible decrease in the objective function from any single variable flip is bounded by $D_{\max}$. Therefore, choosing $\lambda > D_{\max}$ guarantees that the energy cost of violating a constraint always outweighs any potential gain in the objective, forcing the optimizer toward feasible configurations.

\paragraph{Correctness Argument.}
Minimizing $H(\mathbf{x})$ over $\mathbf{x} \in \{0, 1\}^{N(N-1)}$ recovers the optimal TSP tour because the Hamiltonian is constructed such that $H(\mathbf{x}) = O(\mathbf{x})$ for all feasible $\mathbf{x}$ (where all constraints are satisfied and penalties vanish), and $H(\mathbf{x}) > O(\mathbf{x})$ for all infeasible $\mathbf{x}$. Since $\lambda > D_{\max}$, any infeasible configuration has an energy strictly greater than the best feasible configuration. Consequently, the global minimum of $H(\mathbf{x})$ must occur at a feasible point, and among feasible points, $H(\mathbf{x})$ reduces exactly to the objective function $O(\mathbf{x})$. Thus, the global minimizer corresponds precisely to the minimum-distance Hamiltonian cycle.

\paragraph{Illustrative Example.}
Consider a concrete instance with $N = 3$ cities, $V = \{0, 1, 2\}$. The directed edge variables are $x_{01}, x_{02}, x_{10}, x_{12}, x_{20}, x_{21}$. The objective is:
\begin{align}
    O(\mathbf{x}) = d_{01}x_{01} + d_{02}x_{02} + d_{10}x_{10} + d_{12}x_{12} + d_{20}x_{20} + d_{21}x_{21}.
\end{align}
The out-degree and in-degree penalties expand to:
\begin{align}
    P^{\text{out}} &= \lambda \left[ (1 - x_{01} - x_{02} + x_{01}x_{02}) + (1 - x_{10} - x_{12} + x_{10}x_{12}) + (1 - x_{20} - x_{21} + x_{20}x_{21}) \right], \\
    P^{\text{in}} &= \lambda \left[ (1 - x_{10} - x_{20} + x_{10}x_{20}) + (1 - x_{01} - x_{21} + x_{01}x_{21}) + (1 - x_{02} - x_{12} + x_{02}x_{12}) \right].
\end{align}
For $N=3$, the subtour elimination constraint applies only to $S=\{0,1,2\}$, yielding a single penalty term:
\begin{align}
    P^{\text{sub}} = \lambda \left( x_{01} + x_{02} + x_{10} + x_{12} + x_{20} + x_{21} - 2 \right)^2.
\end{align}
Combining all terms, the concrete Hamiltonian is:
\begin{align}
    H(\mathbf{x}) &= \sum_{i \neq j} (d_{ij} - \lambda) x_{ij} + 2\lambda \sum_{\text{quadratic pairs}} x_{ij}x_{kl} + 6\lambda.
\end{align}
For a symmetric distance matrix with $d_{ij} = 1$ for all $i \neq j$ and $\lambda = 2$, the optimal tour is $0 \to 1 \to 2 \to 0$ (or its reverse), corresponding to $x_{01}=x_{12}=x_{20}=1$ and all other $x_{ij}=0$. The objective value is $3$, the penalty terms evaluate to $0$, and the total energy is $H^* = 3$. Any infeasible assignment (e.g., $x_{01}=x_{10}=1$, others $0$) incurs a penalty of at least $2\lambda = 4$, resulting in $H \ge 4 + 2 = 6$, confirming that the global minimum correctly identifies the valid tour.

\end{document}